% ----------------------------------------------------------
% Revisão da Literatura
% ----------------------------------------------------------
\chapter{Revisão da Literatura} \label{cha:revisao}

Este capítulo tem como objetivo apresentar um panorama geral das contribuições científicas à área de orientação vocacional.

\section{Orientação Vocacional}

A trajetória da Orientação Profissional (OP) no Brasil é marcada por transformações teóricas e metodológicas que acompanham as mudanças políticas, sociais e econômicas do país \cite{abade2005}. Historicamente, a produção científica e as práticas na área podem ser compreendidas por meio de três perspectivas principais: a psicométrica, a clínica e a psicossocial \cite{abade2005}.

As primeiras aplicações sistemáticas da Psicologia ao trabalho no Brasil ocorreram na década de 1920, inicialmente voltadas para a seleção profissional em cursos técnicos e empresas ferroviárias \cite{abade2005}. Nesse período, a OP fundamentava-se na administração científica do trabalho, operando sob a lógica do "homem certo no lugar certo" \cite{abade2005}.

O marco institucional dessa fase foi a criação do Instituto de Seleção e Orientação Profissional (ISOP) em 1947, que visava ajustar o trabalhador ao emprego por meio do estudo científico de aptidões \cite{abade2005}. A abordagem era estritamente psicométrica e estatística, utilizando testes vocacionais para direcionar indivíduos de acordo com capacidades inerentes, frequentemente ignorando condições de classe ou histórias de vida \cite{abade2005}.

A partir da década de 1960, a hegemonia da psicometria começou a ser questionada. Sob influência de Carl Rogers e da Psicanálise, o foco deslocou-se do diagnóstico puramente técnico para o auxílio ao autoconhecimento \cite{abade2005}. No entanto, o regime militar instaurado em 1964 restringiu visões críticas, mantendo a Psicologia vinculada a perspectivas experimentais por um longo período \cite{abade2005}.

A consolidação de uma estratégia efetivamente clínica ocorreu no início dos anos 1970 e 1980, fortemente influenciada pela obra do argentino Rodolfo Bohoslavsky \cite{abade2005}. Autores como Maria Margarida Carvalho, na USP, foram pioneiros ao introduzir a estratégia clínica e desenvolver o trabalho de orientação em grupo, que visava ensinar o indivíduo a escolher, integrando aprendizagem experiencial e cognitiva \cite{abade2005}. Nesse período, o termo "vocacional" passou a ser preferido por remeter à dimensão clínica da escolha \cite{abade2005}.

Nos anos 1980 e 1990, surgiram críticas contundentes à OP que ignorava as desigualdades sociais e colocava no indivíduo a responsabilidade pelo sucesso ou insucesso \cite{abade2005}. A escolha profissional passou a ser discutida como um processo multideterminado e dinâmico \cite{abade2005}. Atualmente, a área conta com o suporte da Associação Brasileira de Orientação Profissional (ABOP), criada em 1993, que se tornou o novo centro organizador da produção científica no país \cite{abade2005}. Novas abordagens surgiram, como a abordagem sócio-histórica, que problematiza a realidade social injusta, o paradigma ecológico, fundamentado na Psicologia Social e modelos de ativação da aprendizagem e abordagens integradas baseadas na maturidade vocacional \cite{abade2005}.

A produção científica também destaca as vantagens da OP em grupo. Embora tenha surgido inicialmente para atender a uma demanda quantitativa, percebeu-se que o grupo atua como uma amostra do processo social, permitindo o confronto com a diversidade e o compartilhamento de angústias típicas da adolescência \cite{abade2005}.

Em conclusão, a revisão da literatura indica que, embora as perspectivas clínica e psicométrica ainda predominem, há um movimento crescente em direção a referenciais psicossociais que consideram as transformações do mundo do trabalho contemporâneo e a necessidade de uma prática baseada na realidade social \cite{abade2005}.

\section{Atualidade e Panorama do Mercado}

O setor de tecnologia educacional, ou \textit{EdTech}, atravessa um momento de transição onde a simples digitalização de conteúdos não é mais suficiente; o mercado exige ferramentas que integrem coleta de dados à análise preditiva. Instituições de ensino e profissionais de orientação vocacional lidam hoje com o desafio de escalar o atendimento sem perder a personalização. O \textit{SyncCarreira} atende a essa necessidade crítica ao centralizar o fluxo de dados em uma plataforma única, eliminando a dependência de tabulações manuais que impedem a agilidade no \textit{feedback} ao aluno. Esta automação é o divisor de águas que permite transformar um processo historicamente burocrático em uma estratégia orientada a dados, essencial para a eficácia do aconselhamento profissional moderno.

\section{Aplicação em Outros Contextos}

O \textit{SyncCarreira} foi estruturado com uma arquitetura que permite sua rápida adoção em cenários que exigem alto impacto e baixa complexidade operacional:

\begin{itemize}
    \item \textbf{Escalabilidade em Projetos Sociais e ONGs:} Instituições que operam com jovens em situação de vulnerabilidade sofrem com a limitação de recursos humanos e financeiros para atendimentos individuais. O sistema atua como uma ferramenta de escala: ao automatizar a triagem e o mapeamento de aptidões, o \textit{software} permite que um único psicólogo monitore dezenas de estudantes simultaneamente, garantindo um acompanhamento estruturado que seria inviável via prontuários de papel.

    \item \textbf{Gestão de Grupos em Instituições de Ensino:} A aplicação permite que orientadores gerenciem turmas numerosas sem a fragmentação de informações. Por meio do painel de controle, o profissional tem visibilidade imediata sobre o progresso coletivo, identificando lacunas no desenvolvimento da turma e automatizando o envio de conteúdos direcionados. Isso resolve o gargalo de gestão de grupos, permitindo uma intervenção pedagógica baseada em evidências em vez de suposições manuais.
\end{itemize}
